\begin{table}[htbp]
\centering
\begin{threeparttable}
\caption{Correlations between LLM Readability and Traditional Readability Measures}
\label{tab:readability_corr}
\begin{tabular}{lr}
\toprule
Readability Measure & Pearson $r$ \\
\midrule
Flesch Reading Ease & 0.352 \\
Flesch-Kincaid Grade Level (negated) & 0.460 \\
Gunning Fog Index (negated) & 0.462 \\
SMOG Index (negated) & 0.446 \\
Dale-Chall Score (negated) & 0.273 \\
\bottomrule
\end{tabular}
\begin{tablenotes}
\footnotesize
\item \textit{Notes}. Pearson correlation coefficients between the LLM Readability score (1--10; 10\,=\,easiest to understand) and five traditional readability measures (N\,=\,9,117; AER, ECA, JPE, QJE; English-language articles only; one observation per article). Hengel measures are sign-flipped where needed so that higher values indicate easier/clearer writing: Flesch-Kincaid Grade Level, Gunning Fog Index, SMOG Index, and Dale-Chall Score are negated; Flesch Reading Ease is not.
\end{tablenotes}
\end{threeparttable}
\end{table}
